\documentclass[a4paper]{report}
\usepackage[utf8]{inputenc}
\usepackage[portuguese]{babel}
\usepackage{hyperref}
\usepackage{a4wide}
\hypersetup{pdftitle={CC - TP02},
pdfauthor={João Teixeira, José Ferreira, Miguel Solino},
colorlinks=true,
urlcolor=blue,
linkcolor=black}
\usepackage{subcaption}
\usepackage[cache=false]{minted}
\usepackage{listings}
\usepackage{booktabs}
\usepackage{multirow}
\usepackage{appendix}
\usepackage{tikz}
\usepackage{authblk}
\usepackage{bashful}
\usepackage{verbatim}
\usepackage{amsmath}
\usepackage{tikz}
\usepackage{tikz,fullpage}
\usepackage{pgfgantt}
\usetikzlibrary{arrows,%
                petri,%
                topaths}%
\usepackage{tkz-berge}
\usepackage{titlesec}
\setcounter{secnumdepth}{4}

\titleformat{\paragraph}
{\normalfont\normalsize\bfseries}{\theparagraph}{1em}{}
\titlespacing*{\paragraph}
{0pt}{3.25ex plus 1ex minus .2ex}{1.5ex plus .2ex}
\usetikzlibrary{positioning,automata,decorations.markings}
\AfterEndEnvironment{figure}{\noindent\ignorespaces}
\AfterEndEnvironment{table}{\noindent\ignorespaces}

\begin{document}

\title{Comunicação por Computadores\\ 
\large Fase 2 - Grupo 7}
\author{José Ferreira (A83683) \and João Teixeira (A85504) \and Miguel Solino (A86435)}
\date{\today}

\begin{center}
    \begin{minipage}{0.75\linewidth}
        \centering
        \includegraphics[width=0.4\textwidth]{images/eng.jpeg}\par\vspace{1cm}
        \vspace{1.5cm}
        \href{https://www.uminho.pt/PT}
        {\color{black}{\scshape\LARGE Universidade do Minho}} \par
        \vspace{1cm}
        \href{https://www.di.uminho.pt/}
        {\color{black}{\scshape\Large Departamento de Informática}} \par
        \vspace{1.5cm}
        \maketitle
    \end{minipage}
\end{center}

\pagebreak

\begin{abstract}
    \begin{center}
        Este projeto tem como finalidade implementar uma rede de anonimização de
        tráfego TCP, com recurso a um túnel UDP entre os diferentes nós da rede.
    \end{center}
\end{abstract}

\tableofcontents

\chapter{Introdução}

\chapter{Arquitetura da Solução}

A nossa solução é dividida em duas grandes partes, a parte da comunicação
por TCP com o cliente e o servidor, e a parte do encaminhamento dos pacotes
UDP, interligadas entre um buffer comum a ambas as partes.

Por cada cliente que se liga ao Gateway por TCP, são criadas duas threads, 
uma para lidar com a leitura e outra para lidar com a escrita, de forma
a permitir o tratamento de ambos os fluxos de forma assíncrona. Também no
Gateway de destino são criadas duas threads para a mesma finalidade, mas 
desta vez, para lidar com os fluxos de leitura e escrita do servidor.

Para lidar com os pacotes provenientes de outros Gateway pelo túnel UDP,
é também criada uma thread onde os PDUs são tratados e encaminhados
para a respetiva socket TCP.

\chapter{Especificações do protocolo UDP}
\section{Formato de Mensagens Protocolares (PDU)}

Cada PDU é composto por oito campos, sete para o header e um ultimo
para o conteúdo a transmitir. Um PDU tem um tamanho máximo de 4096,
e, o conteúdo de cada um é, no máximo 1792 bytes, para dar espaço
para a encriptação do mesmo. Caso o conteúdo seja de tamanho 
superior ao que um PDU é capaz de transmitir, este será fragmentado
e identificado de forma a que seja possível a junção de todos os
fragmentos.

\subsection{id}

Este campo identifica sequencialmente cada leitura de uma socket TCP,
tanto originados pelo cliente como pelo servidor de destino. Também 
tem a finalidade de sincronização nas threads de escrita para uma
socket TCP, garantindo assim que pacotes são entregues de forma ordenada.

\subsection{total\_blocks}

Tendo em conta a possível fragmentação do conteúdo por múltiplos PDUs,
este campo tem a informação do número total de PDUs criados para alojar
todo o conteúdo desejado.

\subsection{block\_number}

Aquando da fragmentação do conteúdo por vários PDUs, é atribuído sequencialmente
um identificador de cada fragmento, começando em 0, até \textit{total\_blocks
- 1}, para ser possível a reconstrução do conteúdo fragmentado, e assegurar
a ordem do mesmo.

\subsection{origin\_id}

Para ser possível diferenciar os pacotes e convenientemente entregues ao 
destino, ao serem recebidos pelo túnel UDP, é guardada no PDU o IP do 
Gateway que o enviou.

\subsection{thread\_id}

Para ser possível diferenciar os PDUs provenientes do mesmo Gateway,
é armazenado o identificador único da thread que enviou o PDU.

\subsection{signature}

De forma a ser possível verificar se o PDU foi adulterado ou proveniente
de uma origem que não a desejada, este campo contém a hash do conteúdo
encriptada com a chave privada do Gateway que enviou o PDU.

\subsection{kind} \label{subs:kind}

Para diferenciar os diferentes tipos de PDU, cada PDU tem um campo
com a informação de a que tipo corresponde, com um dos tipos abaixo
listados, para poder ser tratado da forma correta pelo Gateway de 
destino.
\begin{itemize}
        \item Syn
        \item SynAck
        \item Content
        \item Reply
        \item Ack
        \item Fin
        \item FinAck
\end{itemize}

Os pacotes de tipo \textit{Syn} marcam o início de uma conexão, enviado
pelo Gateway ao qual o cliente se ligou por TCP. Os \textit{SynAck}
é a resposta do Gateway que se vai ligar por TCP ao servidor de destino
a um PDU do tipo \textit{Syn}.

Um PDU do tipo \textit{Content} é responsável pela transmissão do conteúdo
de um request TCP proveniente do cliente. Estes PDUs apenas têm origem
no Gateway onde o cliente TCP se conectou. Um PDU do tipo \textit{Reply}
é um PDU análogo ao um do tipo \textit{Content}, mas este apenas tem origem
no Gateway conectado ao servidor de destino.

Os PDU do tipo \textit{Ack} são transmitido com o campo do \textit{content}
vazio, e é enviado como confirmação da receção de outros PDUs.

Os PDU do tipo \textit{Fin} e \textit{FinAck} marcam o fim da transmissão,
e são enviados quando a socket TCP do cliente ou do servidor são fechadas.
Os do tipo \textit{Fin} são enviados pelo Gateway inicial quando é o 
cliente a fechar a conexão, e os \textit{FinAck} são enviados pelo Gateway
conectados ao servidor de destino, quando este fecha a conexão TCP.

\subsection{content}

Este campo é onde é armazenado o conteúdo (ou partes dele) a transmitir pelo
túnel UDP, e tem limite de 1792 bytes por PDU, caso seja um pacote encriptado.

\section{Protocolo de comunicação entre Gateways}

Todas as interaçoes entre os diferentes gateways são efetuadas atraves do túnel
UDP, e estas interaçoes variam conforme o tipo de PDU, descritos em 
\ref{subs:kind}.

\subsection{Protocolo de Início conexão}

Para se establecer uma conexão entre dois nós, o nó no qual o cliente
faz a ligação TCP envia para o nó escolhido um PDU do tipo \textit{Syn},
com a sua chave pública no conteúdo e fica a aguardar o PDU de resposta 
do tipo \textit{SynAck}.

O nó de destino, ao receber um PDU do tipo \textit{Syn}, este guarda
a chave pública da origem, abre a conexão TCP ao servidor de destino, e
responde à origem com o PDU de resposta do tipo \textit{SynAck}, com a sua
chave pública no conteudo.

Ao receber a resposta do nó de destino, a origem guarda também a chave
pública do destino, e dá-se o início da troca de PDUs vindos das 
extremidades TCP.

Estes são os únicos PDUs que são enviados sem autenticação, todos os
posteriores são encriptados e assinados como explicado em \ref{sec:auth}.

\subsection{Protocolo de transmissão de mensagens TCP}

A transmissão das mensagens TCP dá-se por PDUs do tipo \textit{Content} ou
\textit{Reply}. PDUs que têm origem o cliente TCP e destino o servidor,
têm o tipo \textit{Content}, PDUs que fazem o caminho inverso têm o tipo
\textit{Reply}.

PDUs deste tipo são encriptados e assinados antes de serem enviados, e no
nó recetor vão ser validados e desencriptados.

Depois de eviados estes PDUs, a origem fica à espera dos pacotes de confirmação
do tipo \textit{Ack}, e caso estes não cheguem, os PDUs que ainda não receberam
confirmação serão reenviados.

Caso algum PDU não passe a fase de validação, deve ser descartado e não deve
ser enviada a confirmação de receção.

\subsection{Protocolo de finaliização de conexão}

A finaliização da conexão entre dois nós é efetuada atraves de PDUs do tipo
\textit{Fin}, caso seja o cliente a fechar a conexão, ou \textit{FinAck},
caso seja o servidor.

Aquando da receção de um PDU destes tipo, o Gateway deve fechar a conexão TCP
respeitante à origem, e enviar um PDU de confirmação da receção deste, também
do tipo \textit{Ack}.

\chapter{Implementação}

\section{Linguagem de Programação e Bibliotecas utilizadas}

Como linguagem de programação optamos por \textit{Rust}, pela sua sintaxe
apelativa e descomplicada, uma biblioteca standard bastante completa e pela
grande quantidade e qualidade de bibliotecas externas disponíveis. Também
os mecanismos de segurança em acesso a memória partilhada entre threads,
que acontece bastante na nossa solução do projeto, foram um fator que nos
levaram a escolher esta linguagem.

Para além da biblioteca standard de networking da linguagem, com o TCP
listener, stream e UDP sockets, utilizamos também a biblioteca de 
\textit{OpenSSL}, para lidar com o algoritmo \textit{RSA}, e a biblioteca
\textit{serd} e \textit{bincode}, para serialização dos PDUs para binário
e transmissão no túnel UDP, e para a operação inversa no destino. Utilizamos
também uma biblioteca \textit{structopt}, com a finalidade de ler e interpretar 
os argumentos da linha de comandos.

\section{Multiplexagem de clientes e escrita/leitura assíncrona}

De forma a conseguir atender vários clientes ao mesmo tempo por um Gateway,
existe uma thread sempre a correr com o TCP Listener, e cada vez que um cliente
se conecta ao Gateway são criadas duas threads, uma para leitura e uma para
escrita, permitindo assim uma leitura e escrita assíncrona e independente uma
da outra. 

Para garantir o correto encaminhamento das mensagens recebidas no 
túnel UDP, é atribuído um ID a cada cliente, utilizado para ir buscar
as mensagens que a esta lhe dizem respeito, evitando diferentes clientes
receberem mensagens trocadas.

De modo análogo, no Gateway de destino, são criadas também duas threads,
uma para leitura e outra para escrita, indexadas com o mesmo ID do Gateway 
onde o cliente se ligou inicialmente, e estas estão ligadas ao servidor 
de destino.

\section{Ponte TCP-UDP} \label{sec:tcp}

Para fazer a ponte entre o túnel UDP e a socket TCP de destino, ao receber
um PDU pela socket TCP, esta é armazenada num buffer comum a todas as threads
que tratam da escrita para a socket TCP.

O buffer está implementado como sendo um \textit{HashMap}, sendo as chaves um 
par do IP de origem do PDU e o ID da thread TCP da qual foi lido o conteúdo 
correspondente aos PDUs transmitidos. Para cada chave é associado um par que
contem um vetor onde serão armazenados os PDUs, e uma variável de condição,
para a thread responsável por ler do buffer e enviar para a socket TCP correta
ser notificada da existência de novos PDUs. A socket TCP irá ficar em espera
até todos os PDUs da mensagem TCP estarem reunidos no buffer.

\subsection{Garantia de ordem na entrega das mensagens TCP ao destino}

De forma a ser possível garantir que as mensagens TCP chegam de forma ordenada
ao destino, cada thread responsável pela leitura do TCP stream tem um contador,
incrementado sequencialmente a cada mensagem recebida, e anexado aos PDUs
correspondentes. Cada thread responsável pela escrita para o TCP stream, tem
também um contador incrementado sequencialmente a cada mensagem escrita para o
stream. A thread de escrita para o stream apenas procura pelos PDUs com o ID
correspondente ao que se encontra no seu contador. Assim que todo os PDUs com 
esse ID estejam presentes no buffer, os PDUs são ordenados e o seu conteúdo
junto novamente para escrever para o stream.

Como todos os pacotes recebidos de um lado têm que ser entregues do outro,
os contadores das diferentes threads estarão sempre sincronizados, assegurando
assim a ordem de entrega.

\section{Controlo de perdas}

Para evitar as perdas de PDUs e garantir que todos chegam ao destino, a cada 
PDU recebido no túnel UDP, é enviado em reposta a este um outro com a
informação de que chegou ao destino e é válido. O Gateway que envia o PDU
fica à espera que estes PDUs de confirmação cheguem, durante um determinado
intervalo de tempo, que recomeça cada vez que um PDU de confirmação chega.
Caso o tempo passe sem chegar nenhuma nova confirmação, os pacotes que ainda
não chegou a confirmação de receção são retransmitidos.

Os PDU de confirmação são encaminhados para um buffer diferente, com o mesmo
tipo e funcionamento do descrito em \ref{sec:tcp}.

\section{Confidencialidade e Assinatura digital} \label{sec:auth}

Para comprovar que o PDU é proveniente da origem correta, que não foi
adulterado e que o conteúdo se manteve confidencial durante a transmissão,
encriptamos e assinamos o PDU com recurso algoritmo de chaves públicas 
\textit{RSA}. 

Para assegurar que a origem do PDU, é efetuada o \textit{hashing} do conteúdo
com algoritmo \textit{SHA-2}, na variante de 512 bits, e em seguida, encriptada
com a chave privada do Gateway de origem, garantindo assim a autenticação da
origem, o não repudio e a integridade do conteúdo. Para efetuar esta validação,
no gateway de destino, é feita a desencriptação da hash, e a comparação desta 
com a hash do conteúdo contido. Caso as duas hashs não sejam iguais, o PDU é
descartado.

Para assegurar a confidencialidade do conteúdo do PDU, este é encriptado com
recurso à chave pública do Gateway de destino, assegurando assim que este
apenas pode ser desencriptado pelo correto destinatário. O conteúdo é 
encriptado em blocos de 128 bytes, e estes blocos após encriptados, são
concatenados até ao tamanho máximo de um PDU e encaminhados para o túnel
UDP.

\chapter{Testes e Resultados}

\chapter{Conclusões e Trabalho Futuro}

\end{document}
